\documentclass{article}
\usepackage{amsmath,amsfonts,amssymb,enumitem}
\usepackage[english]{babel}

\newcommand{\openquote}{\textquotedblleft}
\newcommand{\closequote}{\textquotedblright}
\usepackage{listings}

\title{MAT 211: Homework \#4}
\date{Spring 2026}
\author{Your Name}

\begin{document}
\maketitle 

\lstset{
  basicstyle=\ttfamily\small,
  frame=single,
  breaklines=true
}

\begin{enumerate}

\item The following code uses the \emph{rel} tactic in Lean.

\begin{lstlisting}

example {r s : Q} (h1 : s + 3 >= r) (h2 : s + r <= 3) : r <= 3 :=
  calc
    r = (s + r + r - s) / 2     := by ring
    _ <= (3 + (s + 3) - s) / 2  := by rel[h1,h2]
    _ = 3                       := by ring

\end{lstlisting}

Prove the following \begin{quote}
    If $r, s \in \mathbb{Q}$ satisfy
    $s + 3 \geq r$ and $s + r \leq 3$, then $r \leq 3$
  \end{quote} without appealing to Lean tactics.
  That is, write a clear mathematical argument that the conclusion $r \leq 3$ follows from the
  two hypotheses. Your proof should reflect the logical structure of the Lean \texttt{calc}
  block above.
 

    \item Consider the following two Lean proofs of the same inequality.
The first succeeds; the second produces a Lean error on the \texttt{rel} step.

\begin{lstlisting}[caption={Proof A (succeeds)}]
example {t : Q} (ht : t >= 10) : t ^ 2 - 3 * t - 17 >= 5 :=
  calc
    t ^ 2 - 3 * t - 17
        = t * t - 3 * t - 17       := by ring
      _ >= 10 * t - 3 * t - 17     := by rel[ht]
      _ = 7 * t - 17               := by ring
      _ >= 7 * 10 - 17             := by rel[ht]
      _ >= 5                       := by numbers
\end{lstlisting}

\newpage 

\begin{lstlisting}[caption={Proof B (fails)}]
example {t : Q} (ht : t >= 10) : t ^ 2 - 3 * t - 17 >= 5 :=
  calc
   t ^ 2 - 3 * t - 17
    >= 10 ^ 2 - 3 * 10 - 17    := by rel[ht]   -- ERROR
    _ = 87                     := by numbers
    _ >= 5                     := by numbers
\end{lstlisting}

\medskip
\noindent
After comparing these proofs, answer the following questions.

\begin{enumerate}
    \item The \texttt{rel} tactic checks whether a substitution is valid by inspecting the \emph{syntactic structure} 
    of the expression---it walks the expression tree and verifies, node by node, that replacing $t$ with a larger value moves each sub-expression in a
    \emph{consistent} direction.  We call this property \textbf{syntactic monotonicity}.\\

    In Proof~B, the \texttt{rel} step asks Lean to replace $t$ with $10$
    throughout $t^2 - 3t - 17$ in a single step.
    Identify the two sub-expressions of $t^2 - 3t - 17$ that
    \emph{conflict} with each other under this substitution, and explain
    in one or two sentences why their combined effect prevents
    \texttt{rel} from certifying the inequality.

\medskip
\item 
Proof~A avoids this conflict by splitting the argument into steps
where each \texttt{rel} application is unambiguous.
Examine the first \texttt{rel} step in Proof~A:
\[
  t \cdot t - 3t - 17 \;\geq\; 10 \cdot t - 3t - 17.
\]
Only \emph{one} factor of the product $t \cdot t$ is replaced by
$10$, while the $-3t$ term is left unchanged on both sides.
Explain why this makes the monotonicity check straightforward for
\texttt{rel}.  (\emph{Hint:} what sign does the remaining factor of
$t$ have, and why does that matter?)

\medskip
\item 
After the first \texttt{rel} step, a \texttt{ring} step simplifies
$10t - 3t - 17$ to $7t - 17$, and then a second \texttt{rel} step
handles the remaining inequality $7t - 17 \geq 70 - 17$.
Why is this second \texttt{rel} step unambiguous, even though $t$
still appears in a subtraction?

\item  (\emph{Reflection.})
A more experienced Lean user might instead close the goal in one
line with \texttt{nlinarith} or a similar tactic.
Give one reason why the step-by-step \texttt{calc} proof in Proof~A
might be \emph{preferable} as a learning exercise, even if it is
longer.
\end{enumerate}

\newpage 

\textbf{Prove either directly, by contrapositive, by contradiction as appropriate.}

\item Prove for positive real numbers $a$ and $b$, \(a+b \ge 2 \sqrt{ab}\).

\item Show for \(x\in\mathbb{Z}\), $x$ is odd if and only if $3x+6$ is odd.

\item Prove \(\forall x\in\mathbb{Z}\), \(x^3\) is even if and only if \(x\) is even.

\item Prove \( \sqrt[3]{2} \) is irrational. You may use the previous exercise as a lemma, if needed.
    
\item Prove for \(x,y\in\mathbb{R^+}, \sqrt{x+y}\ne\sqrt{x} +\sqrt{y}\).  Note: \(\mathbb{R^+}\) represents the positive real numbers.

\item Prove: If $p$ is prime and $p\mid a^2$, then $p\mid a$.

\item Prove: If $a \in\mathbb{Z}$, then $a^3 \equiv a$ (mod 3).

\end{enumerate}

\end{document}